\section{Shifted ranks and successive norm transport: a second parallel continuation}
\label{sec:rank-transport-continuation}
This continuation follows different unresolved gates from the preceding round.
On the forward arithmetic route, exact shifted residue classes give an
elementary average bound for the full multiplicity of primes in a growing
range. Retaining the residual norm in the height denominator makes the bound
uniform in the residual and in the root. Finite unions then permit adaptive
selection from specified growing norm boxes. Exceptional indices and primes
outside the controlled range remain explicit.

A further two-place refinement depends on the remainder proportion
$\lambda=\max(1,\log N(v))/g$ alone. It treats the ramified sign, unit
residuals and multiplicatively dependent factors explicitly. Actual
representations with $\lambda\to0$ and displayed $\rho\to\infty$ demonstrate
the gain over the earlier sufficient test. The refinement also proves a
positive second-$\lambda$ floor when the first root is fixed and its first
$\lambda$ is bounded. It does not compare that floor with the finite
two-step threshold or bound it uniformly for moving first roots.

On the reverse route, actual radical transport through two norm transforms
gives a sufficient criterion for a disproof. A fixed quartic field and unit
describe the second norm's prime depths. Local lifting supplies simultaneous
arbitrarily deep divisibility at two distinct primes on actual primitive
families. Such local compatibility does not give compression of the entire
second norm, and the required unbounded family is still unconstructed.

The homogeneous signed-rank decomposition remains a separate route. Its
rankwise polynomial hypothesis implies that the positive part of the normalized
signed tail tends to zero. An upper bound tending to zero is not asserted to
be a two-sided limit for the signed quantity itself. The residual counterfamily
refutes a precise homogeneous lifting rule transferred to shifted orbits;
its fixed prime eventually lies below the moving tail cutoff.

\paragraph{Verification boundary.}
The ordinary proofs in the next sections were checked by other research agents
before integration. The new formal modules cover actual residue-class
parameterizations and exact counting errors, an infinite finite-depth orbit,
integer radical identities and compression, and an implication from an
explicit unbounded-family premise to the negation of the repository's standard
$abc$ definition. The family premise is not supplied. The analytic estimates,
complete local group classification, and full ordinary research program are
not silently included in that formal scope. Independent agent review does not
constitute external peer review.

\paragraph{Continuing obligations.}
The forward route still needs a signed large-prime bound and control of
exceptional indices. The reverse route still needs either an actual unbounded
family meeting all compression conditions or a theorem excluding a precisely
stated family. Moving roots, shared remainders, and the broader geometric,
IUT, Pell, Mersenne and incidence routes stay open unless their full premises
are decided by proof or counterexample. This continuation does not prove or
disprove standard $abc$.
